\section{Estado del arte}

\subsection{Enfoques generales de PLN para detección de salud mental}

Los trabajos de revisión recientes muestran que la detección de salud mental mediante PLN ha evolucionado desde representaciones manuales basadas en palabras, léxicos y metadatos hacia arquitecturas profundas y modelos basados en transformers, pero con una persistencia notable de problemas estructurales. \citet{montejo2024survey} sintetizan tendencias, datasets, algoritmos y campañas de evaluación en varios trastornos, y subrayan carencias recurrentes: predominio del inglés, necesidad de anotación experta, bajo énfasis en explicabilidad y escasa validación en escenarios reales. Desde una perspectiva más específica, \citet{abdulsalam2024suicidal} revisan treinta y siete estudios sobre ideación suicida y confirman que el campo sigue concentrado en Twitter, Reddit y Weibo, con una mezcla de aprendizaje supervisado clásico, CNN, LSTM y modelos híbridos.

Estas revisiones son importantes porque desplazan la discusión desde el rendimiento aislado de cada artículo hacia la madurez metodológica del área. En términos de efectividad predictiva, ambos estudios muestran que el rendimiento depende no solo del modelo, sino también del esquema de anotación, el idioma, la plataforma y la granularidad de la tarea. En términos de interpretabilidad clínica, ambos coinciden en que la literatura tiende a operacionalizar el problema como clasificación textual y no como apoyo explicativo a la toma de decisiones, lo que deja una brecha entre desempeño computacional y utilidad clínica \citep{montejo2024survey, abdulsalam2024suicidal}.

Un antecedente clave para comprender esta tensión es el trabajo que modela el desplazamiento desde discurso general de salud mental hacia discurso suicida en Reddit, en lugar de limitarse a clasificar publicaciones aisladas \citep{dechoudhury2016shifts}. Su contribución es relevante porque introduce una lógica longitudinal y conductual: el riesgo no se expresa únicamente por palabras explícitas, sino por cambios en cohesión lingüística, foco en uno mismo, interacción social y marcadores de desesperanza. Desde la dimensión de interpretabilidad clínica, ese cambio de perspectiva aumenta el valor del análisis, ya que permite pensar en detección temprana y no solo en clasificación retrospectiva. Este trabajo aporta al reto porque muestra que la utilidad del sistema depende también de capturar trayectorias discursivas y no únicamente etiquetas finales.

\subsection{Detección de ideación suicida en redes sociales}

En la línea de clasificación de publicaciones suicidas, \citet{tadesse2020suicide} representan una etapa de transición entre aprendizaje automático convencional y aprendizaje profundo. Su combinación LSTM-CNN con \textit{word embeddings} mejora sobre modelos aislados y sobre clasificadores clásicos, mostrando que las arquitecturas híbridas capturan simultáneamente dependencias secuenciales y patrones locales. Desde la dimensión de efectividad predictiva, el trabajo es importante porque legitima el uso de redes profundas sobre Reddit; sin embargo, su interpretabilidad es reducida y los propios autores reconocen limitaciones de tamaño y sesgo en los datos.

\citet{chatterjee2022multimodal} siguen otra trayectoria: en vez de depender únicamente del texto bruto, construyen una técnica multiatributo que combina rasgos lingüísticos, temáticos, sentimentales, temporales, de personalidad y de comportamiento en Reddit y Twitter. El valor de este enfoque radica en que la mejora predictiva no proviene solamente de una arquitectura más compleja, sino de una representación del problema más rica. Desde la dimensión de efectividad predictiva, esto permite integrar señales complementarias; desde la dimensión de interpretabilidad clínica, introduce variables con mayor potencial explicativo, pues el tiempo de publicación, los emoticonos, los rasgos afectivos y los perfiles de interacción pueden convertirse en señales más legibles que un vector contextual opaco.

La evaluación de severidad propuesta por \citet{gaur2019severity} profundiza todavía más en la conexión entre señal computacional y significado clínico. En lugar de clasificación binaria, los autores formulan un esquema de cinco niveles inspirado en C-SSRS, diferencian usuarios de apoyo de usuarios con diferentes grados de riesgo y utilizan recursos médicos para crear un léxico de severidad suicida. Su CNN logra el mejor equilibrio entre \textit{graded recall}, error ordinal y medida de riesgo percibido. A diferencia de los enfoques puramente predictivos, esta formulación es especialmente valiosa para la interpretabilidad clínica porque organiza la salida del sistema en niveles accionables y no solo en etiquetas abstractas. Este trabajo aporta al reto porque aproxima la clasificación automática a decisiones de triage más cercanas al lenguaje profesional.

Los trabajos más recientes refuerzan la búsqueda de mayor desempeño con arquitecturas profundas. \citet{mirtaheri2024tcn} proponen un modelo LSTM--TCN con autoatención sobre Reddit y Twitter, argumentando que la combinación mejora memoria de largo alcance, escalabilidad y extracción de rasgos relevantes. \citet{shuvo2024bigru} comparan varios modelos de lenguaje preentrenados acoplados con Bi-GRU y reportan un rendimiento muy alto, además de bajas tasas de falsos negativos, aspecto especialmente crítico en aplicaciones de riesgo. Desde la dimensión de efectividad predictiva, ambos artículos refuerzan la frontera cuantitativa del campo. No obstante, ambos muestran una limitación común del estado del arte: aun cuando la precisión aumenta, la explicación clínica de por qué una publicación fue catalogada como riesgosa sigue siendo secundaria.

\subsection{Detección de depresión y enfoques clínicamente fundamentados en DSM-5}

La literatura sobre depresión en redes sociales ha pasado de la detección general del trastorno a esquemas más cercanos a la práctica clínica. \citet{gaur2018dsm5} representan un antecedente importante al mapear subreddits de salud mental a categorías DSM-5 mediante modelado temático, \textit{embeddings} y bases de conocimiento médico. Su propuesta SEDO reduce de manera sustancial la tasa de falsas alarmas y obtiene una concordancia alta con expertos. Desde la dimensión de interpretabilidad clínica, el trabajo es central porque introduce una mediación semántica entre lenguaje social y categorías diagnósticas. Aunque el objeto de análisis son subreddits y no pacientes individuales, el estudio demuestra que la incorporación de conocimiento clínico estructurado puede mejorar simultáneamente precisión y trazabilidad semántica.

Ese movimiento hacia una mayor granularidad clínica se consolida con \citet{bao2025redsm5}, quienes presentan ReDSM5, un corpus de 1484 publicaciones largas de Reddit anotadas a nivel de oración por una psicóloga con los nueve síntomas depresivos del DSM-5 y racionales clínicos breves. Este recurso es uno de los aportes más importantes del corpus revisado para la dimensión de interpretabilidad clínica: la tarea deja de ser ``depresión sí/no'' y pasa a ser identificación multietiqueta de síntomas concretos, con evidencia textual y explicación. Desde la dimensión de efectividad predictiva, los resultados base muestran que los modelos basados en transformers ajustados superan con claridad a SVM y CNN, lo que sugiere que ambas metas pueden coexistir cuando el diseño del dataset lo permite. Este trabajo aporta al reto porque ofrece una referencia directa de cómo vincular señales lingüísticas con síntomas clínicos observables.

\citet{qasim2025depression} complementan esta línea al centrarse en la severidad depresiva y comparar enfoques basados en n-gramas, \textit{sentence transformers} y modelos basados en transformers ajustados de extremo a extremo. Sus resultados muestran que modelos como DeBERTa capturan mejor las diferencias entre clases leve, moderada y severa que las representaciones estáticas o los modelos con ingeniería manual de atributos. Desde la dimensión de efectividad predictiva, el artículo confirma la ventaja de los modelos basados en transformers finamente ajustados. Aun así, mantiene una interpretabilidad limitada en comparación con los trabajos DSM-5, pues la explicación se apoya principalmente en el poder contextual del modelo y no en racionales clínicos explícitos.

\subsection{Aprendizaje profundo, modelos basados en transformers y modelos de lenguaje de gran escala}

La fase más reciente del estado del arte desplaza el foco desde la simple clasificación hacia modelos capaces de representar contexto amplio, extraer evidencia y apoyar interpretación. \citet{bauer2024llm} utilizan \textit{embeddings} de modelos de lenguaje, reducción de dimensionalidad y apoyo de un modelo generativo para analizar 2.9 millones de publicaciones de treinta subreddits. Su objetivo principal no es predecir riesgo individual, sino comprender dimensiones latentes del lenguaje suicida. Desde la dimensión de interpretabilidad clínica, este artículo es especialmente relevante, ya que conecta patrones lingüísticos con teorías contemporáneas del suicidio como desconexión, carga percibida, desesperanza o trauma.

Por su parte, \citet{singh2024evidence} exploran una tarea todavía más cercana al apoyo experto: extraer y resumir evidencia textual de ideación suicida con LLMs. Su desempeño sobresaliente en la tarea compartida de CLPsych 2024 muestra que los modelos generativos pueden producir salidas útiles no solo para clasificar, sino también para justificar. A diferencia de los enfoques puramente predictivos, esta transición sitúa la interpretabilidad en el nivel de fragmentos de evidencia y resúmenes estructurados, no únicamente en variables agregadas o pesos de atención. Este trabajo aporta al reto porque muestra una vía concreta para convertir el resultado del modelo en evidencia comunicable.

En conjunto, los artículos de esta sublínea sugieren una conclusión provisional. Los modelos basados en transformers y los LLMs elevan claramente la capacidad del campo para modelar contexto, severidad y relaciones semánticas complejas; sin embargo, su utilidad clínica depende de cómo se diseñe el problema. Cuando se usan solo para maximizar exactitud, la interpretabilidad sigue siendo débil. Cuando se combinan con anotación experta, taxonomías clínicas o tareas de extracción de evidencia, se convierten en herramientas más pertinentes para salud mental \citep{bao2025redsm5, bauer2024llm, singh2024evidence}.

Esta evolución también muestra que los modelos de lenguaje no sustituyen por completo a los enfoques previos, sino que amplían el tipo de preguntas que pueden formularse. Mientras los modelos clásicos y profundos se concentran principalmente en clasificación, los LLMs permiten explorar tareas de extracción de evidencia, generación de explicaciones y análisis semántico de mayor granularidad. Para el reto, esta diferencia es importante porque sugiere que la comparación futura entre enfoques debe considerar tanto el desempeño numérico como el tipo de salida producida por cada modelo.
